\documentclass[11pt,letterpaper]{article}
\usepackage[margin=1in]{geometry}
\usepackage{booktabs}
\usepackage{tabularx}
\usepackage{longtable}
\usepackage{array}
\usepackage{hyperref}
\usepackage{amsmath}
\usepackage{parskip}
\usepackage{titlesec}
\usepackage{enumitem}
\usepackage[numbers]{natbib}

\titlespacing*{\section}{0pt}{1.2ex plus .2ex}{0.8ex}
\titlespacing*{\subsection}{0pt}{1ex plus .2ex}{0.5ex}

\hypersetup{colorlinks=true, linkcolor=blue, citecolor=blue, urlcolor=blue}

\title{\textbf{AI Virtual Cell Models for Dopaminergic Neuron Perturbation Prediction}\\
\large Midterm Progress Report}
\author{Katherine Deborah Godwin Gnanaraj \and Zihan Jin \and Yumejichi Fujita \and Brandon Latherow}
\date{March 2026 \quad CSCI 566}

\begin{document}
\maketitle

%% ─────────────────────────────────────────────────────────────────────────────
\begin{abstract}
Parkinson's disease (PD) is driven by selective loss of dopaminergic (DA) neurons in the
substantia nigra. AI virtual cell (AIVC) models trained on single-cell CRISPR perturbation
screens can predict post-perturbation gene expression, enabling in silico gene ranking for
experimental validation. We benchmark a mean expression baseline, scGen, GEARS, STATE/CPA,
and scGenePT on the GSE152988 iPSC-derived neuron dataset (CRISPRi and CRISPRa modalities).
All models are evaluated under both absolute Pearson metrics and stricter delta-based metrics
that isolate perturbation-specific signal from cell-type background. GEARS achieves Pearson
$\Delta$ (top-20 DEGs) of 0.262 on completely unseen gene knockdowns; sign accuracy of 79.3\%
on CRISPRa demonstrates meaningful directional prediction on the DA-marker gene panel. Results
from STATE/CPA (Zihan Jin) and scGenePT (Brandon Latherow) are being integrated.
\end{abstract}

%% ─────────────────────────────────────────────────────────────────────────────
\section{Introduction}

CRISPR-based perturbation screens combined with single-cell RNA-seq can measure how hundreds of
gene knockdowns reshape transcriptional profiles at single-cell resolution. Systematically
testing all $\sim$20,000 human protein-coding genes experimentally remains infeasible, making
computational perturbation prediction a key tool for prioritizing candidates. In this project we
benchmark several AIVC models on iPSC-derived neuron datasets to identify which gene
perturbations most strongly promote or inhibit a DA neuron transcriptional state --- a central
question for Parkinson's disease therapeutic target discovery.

%% ─────────────────────────────────────────────────────────────────────────────
\section{Related Work}

\textbf{scGen}~\cite{lopez2018} is a VAE that models perturbation effects via latent space
arithmetic. It is fast and interpretable but cannot generalize to unseen perturbations by design.

\textbf{GEARS}~\cite{roohani2023} adds a gene co-expression graph (STRING) to a GNN, enabling
prediction of gene knockdowns never observed in training. It is the standard baseline for
zero-shot perturbation evaluation.

\textbf{CPA}~\cite{lotfollahi2021} learns disentangled latent representations for cell identity
and perturbation, enabling compositional prediction of perturbation combinations.

\textbf{STATE} is the CZ AIVC portal framework supporting multiple backends (STATE, CPA, scVI,
scGPT, Tahoe). We fine-tune the \texttt{ST-SE-Replogle} pretrained model on GSE152988 via this
framework.

\textbf{scGenePT}~(2024) augments scGen with gene-level text embeddings from biological
literature, improving generalization to unseen perturbations.

\textbf{scFoundation}~\cite{hao2024} is a 50M-cell pretrained foundation model generating
context-aware gene embeddings (19,264-gene vocab) for use as richer GEARS node features.
Planned for evaluation if compute permits.

\begin{table}[h]
\centering
\caption{Models surveyed and evaluated.}
\small
\begin{tabular}{llllll}
\toprule
Model & Architecture & Unseen Perts & Gene Graph & Status \\
\midrule
Mean Baseline        & ---                  & No      & No      & Done \\
scGen                & VAE                  & No      & No      & Done \\
GEARS (5 ep, 5k HVG) & GNN                 & Yes     & STRING  & Done \\
GEARS (20 ep, all genes) & GNN             & Yes     & STRING/None$^*$ & Done \\
STATE (GSE152988 combined) & Structured+CPA/scVI & Partial & No & In progress \\
scGenePT             & VAE + LLM text       & Partial & No      & In progress \\
scFoundation + GEARS & GNN + Transformer    & Yes     & STRING  & Planned \\
\bottomrule
\end{tabular}
\vspace{2pt}
{\footnotesize $^*$One run used an empty graph due to a library compatibility issue with the full gene set.}
\end{table}

%% ─────────────────────────────────────────────────────────────────────────────
\section{Datasets and Preprocessing}

\subsection*{GSE152988 --- Primary Benchmark}

iPSC-derived neurons~\cite{tian2021} with two CRISPR modalities:
\begin{itemize}[noitemsep]
  \item \textbf{CRISPRi:} 32,300 cells, 184 gene knockdowns, 437 controls, 33,538 genes
  \item \textbf{CRISPRa:} 21,193 cells, 100 gene activations, 434 controls, 33,538 genes
\end{itemize}
Raw counts normalized to 10,000 counts/cell and log1p-transformed. 67/71 DA marker genes
present in the panel, confirming suitability for DA neuron transcriptional studies.

\subsection*{GSE124703 --- Secondary / Zero-Shot Transfer Target}

iPSC and Day 7 neuron cells~\cite{tian2019} with 57 CRISPRi guides across two developmental
contexts, enabling zero-shot cell-type transfer evaluation.

\subsection*{Data Split}

\textbf{Within-dataset (GEARS, scGen, scGenePT):} \texttt{zero\_shot\_pert} split --- conditions
with $\geq$30 cells are split 80/10/10 (train/val/test) by condition, seed = 42. Controls always
in training. CRISPRi: 148/18/18; CRISPRa: 80/10/10. Test conditions are fully held out.

\textbf{Cross-dataset (STATE):} Fine-tune on combined GSE152988 (228/28/28 perturbations) using
modality-specific labels (\texttt{CRISPRI::GENE}, \texttt{CRISPRA::GENE}); evaluate zero-shot on
GSE124703.

\textbf{Data leakage note:} For pretrained models (STATE/\texttt{ST-SE-Replogle}, scFoundation),
GSE152988 or GSE124703 may be present in the pretraining corpus. Results should be interpreted
with this caveat until verified against pretraining manifests.

%% ─────────────────────────────────────────────────────────────────────────────
\section{Evaluation Framework}

We use a two-tier metric suite that separates overall expression reconstruction (secondary) from
perturbation-specific delta prediction (primary). Absolute Pearson across all genes is dominated
by cell-type identity and will exceed 0.99 for even a trivial mean baseline, making it
uninformative about perturbation prediction quality.

\textbf{Primary --- delta-based:} For each held-out perturbation, define
$\text{true\_delta} = \bar{x}_{\text{pert}} - \bar{x}_{\text{ctrl}}$ and
$\text{pred\_delta} = \hat{x}_{\text{pert}} - \bar{x}_{\text{ctrl}}$, then compute:
\begin{itemize}[noitemsep]
  \item \textbf{Pearson $\Delta$ top-20 DEGs} --- correlation of predicted vs.\ true \emph{change} on the 20 most differentially expressed genes
  \item \textbf{Pearson $\Delta$ DA markers} --- same restricted to the 71-gene DA marker panel
  \item \textbf{Sign accuracy (top-20 DEGs / DA markers)} --- fraction where predicted up/down direction matches
\end{itemize}

\textbf{Secondary --- absolute:} Pearson $r$ (all genes), Pearson $r$ (top-20 DEGs), Pearson $r$
(DA markers), MSE.

%% ─────────────────────────────────────────────────────────────────────────────
\section{Results}

\subsection{Protocol Comparison}

The model families use different evaluation protocols and \textbf{cannot be ranked directly}:

\begin{table}[h]
\centering
\small
\begin{tabular}{llll}
\toprule
Model & Split Type & \# Test Perts & Task \\
\midrule
Mean Baseline & Seen   & 20    & Predict global mean \\
scGen         & Seen   & 20    & Seen perturbations, 80/20 cell split \\
GEARS         & \textbf{Unseen} & 18--46 & Entire conditions held out \\
\bottomrule
\end{tabular}
\end{table}

\subsection{CRISPRi Results}

\begin{table}[h]
\centering
\caption{CRISPRi --- primary delta-based metrics.}
\small
\begin{tabular}{lcccccc}
\toprule
Model & Pearson $\Delta$ top-20 & Pearson $\Delta$ DA & Sign acc top-20 & Sign acc DA & \# Perts & Protocol \\
\midrule
Mean Baseline         & 0.769   & 0.518 & 0.953 & 0.653 & 20 & Seen \\
scGen                 & 0.442   & 0.168 & 0.738 & 0.478 & 20 & Seen \\
GEARS (5 ep, 5k HVG)  & 0.262$^\dagger$ & --- & 0.663$^\dagger$ & --- & 46 & \textbf{Unseen} \\
GEARS (20 ep)         & 0.129$^\dagger$ & --- & 0.623$^\dagger$ & --- & 37 & \textbf{Unseen} \\
\bottomrule
\end{tabular}
\vspace{2pt}
{\footnotesize $^\dagger$From GEARS' internal evaluation. DA marker delta omitted for CRISPRi: small effect sizes make delta sensitive to ctrl\_mean reference; absolute DA metrics (Table~3) are reliable.}
\end{table}

\begin{table}[h]
\centering
\caption{CRISPRi --- secondary absolute metrics.}
\small
\begin{tabular}{lccccc}
\toprule
Model & Pearson $r$ (all) & Pearson $r$ (top-20) & Pearson $r$ (DA) & MSE & Protocol \\
\midrule
Mean Baseline        & 0.9988 & 0.9917 & 0.9998 & 0.0005  & Seen (20) \\
scGen                & 0.9938 & 0.9523 & 0.9988 & 0.0025  & Seen (20) \\
GEARS (5 ep)         & 0.9965 & 0.9733 & ---    & 0.00137 & \textbf{Unseen (46)} \\
GEARS (20 ep)        & 0.9963 & \textbf{0.9776} & 0.9984 & 0.0007 & \textbf{Unseen (37)} \\
\bottomrule
\end{tabular}
\end{table}

All models achieve Pearson $r > 0.99$ across all genes including the mean baseline, confirming
this metric reflects cell-type identity rather than perturbation prediction.
The mean baseline's high delta Pearson (0.769) is an artifact: randomly sampled seen perturbations
are representative of the dataset average, consistent with findings that simple baselines are
competitive on seen perturbations~\cite{boiarsky2023,wenk2024}. GEARS' sign accuracy of
\textbf{62.3\%} on completely unseen perturbations is above the 50\% chance baseline.

\subsection{CRISPRa Results}

\begin{table}[h]
\centering
\caption{CRISPRa --- primary delta-based metrics.}
\small
\begin{tabular}{lcccccc}
\toprule
Model & Pearson $\Delta$ top-20 & Pearson $\Delta$ DA & Sign acc top-20 & Sign acc DA & \# Perts & Protocol \\
\midrule
Mean Baseline        & 0.591  & 0.628 & 0.920 & 0.493 & 20 & Seen \\
scGen                & 0.591  & 0.283 & 0.833 & 0.407 & 20 & Seen \\
GEARS (5 ep)         & 0.334$^\dagger$ & --- & 0.714$^\dagger$ & --- & 25 & \textbf{Unseen} \\
GEARS (20 ep)        & \textbf{0.398} & \textbf{0.498} & \textbf{0.793} & 0.510 & 20 & \textbf{Unseen} \\
\bottomrule
\end{tabular}
\end{table}

\begin{table}[h]
\centering
\caption{CRISPRa --- secondary absolute metrics.}
\small
\begin{tabular}{lccccc}
\toprule
Model & Pearson $r$ (all) & Pearson $r$ (top-20) & Pearson $r$ (DA) & MSE & Protocol \\
\midrule
Mean Baseline        & 0.9956 & 0.9738 & 0.9964 & 0.0018  & Seen (20) \\
scGen                & 0.9928 & 0.9582 & 0.9960 & 0.0034  & Seen (20) \\
GEARS (5 ep)         & 0.9948 & 0.9565 & ---    & 0.00224 & \textbf{Unseen (25)} \\
GEARS (20 ep)        & 0.9941 & \textbf{0.9614} & 0.9937 & 0.0010 & \textbf{Unseen (20)} \\
\bottomrule
\end{tabular}
\end{table}

CRISPRa shows substantially stronger GEARS performance. Pearson $\Delta$ top-20 DEGs of
\textbf{0.398} and DA marker delta Pearson of \textbf{0.498} at 20 epochs, with sign accuracy
of \textbf{79.3\%} on completely unseen gene activations. The DA marker Pearson delta of 0.498
is the most relevant result for this project: the model correctly predicts the direction and
magnitude of perturbation effects on canonical DA marker genes (TH, SLC6A3, NR4A2, FOXA2) for
unseen perturbations. Gene activation (CRISPRa) produces larger transcriptional responses than
repression, making delta metrics more reliable.

\subsection{Additional Model Results}

\textit{[To be added by teammates: STATE/CPA results (Zihan Jin), scGenePT results (Brandon Latherow)]}

%% ─────────────────────────────────────────────────────────────────────────────
\section{Roadmap}

\begin{itemize}[noitemsep]
  \item \textbf{Immediate:} Integrate STATE/CPA and scGenePT results into Tables 2--5
  \item \textbf{Immediate:} Zero-shot cross-dataset evaluation --- STATE fine-tuned on GSE152988 $\to$ test on GSE124703 (unseen developmental context)
  \item \textbf{If compute permits:} scFoundation + GEARS on HPC or Colab A100 (requires 40+ GB VRAM)
  \item \textbf{Final report:} DA identity scoring --- rank perturbations by predicted shift toward DA neuron transcriptional state using GSE140231 reference
\end{itemize}

%% ─────────────────────────────────────────────────────────────────────────────
\bibliographystyle{unsrtnat}
\begin{thebibliography}{9}

\bibitem{lopez2018}
Lopez, R. et al.\ (2018).
Deep generative modeling for single-cell transcriptomics.
\textit{Nature Methods}, 15, 1053--1058.

\bibitem{roohani2023}
Roohani, Y. et al.\ (2023).
Predicting transcriptional outcomes of novel multigene perturbations with GEARS.
\textit{Nature Biotechnology}, 42, 927--935.

\bibitem{hao2024}
Hao, M. et al.\ (2024).
Large-scale foundation model on single-cell transcriptomics.
\textit{Nature Methods}, 21, 1481--1491.

\bibitem{lotfollahi2021}
Lotfollahi, M. et al.\ (2021).
Compositional perturbation autoencoder for single-cell response modeling.
\textit{bioRxiv}.

\bibitem{tian2021}
Tian, R. \& Kampmann, M. (2021).
Genome-wide CRISPRi/a screens in human neurons link lysosomal failure to ferroptosis.
\textit{Nature Neuroscience}, 24, 1020--1034.

\bibitem{tian2019}
Tian, R. et al.\ (2019).
CRISPR interference-based platform for multimodal genetic screens in human iPSC-derived neurons.
\textit{Neuron}, 104, 239--255.

\bibitem{boiarsky2023}
Boiarsky, R. et al.\ (2023).
A single cell foundation model.
\textit{bioRxiv}.

\bibitem{wenk2024}
Wenk, P. et al.\ (2024).
Simple baselines for perturbation modelling in single-cell data.
\textit{bioRxiv}.

\end{thebibliography}

\vspace{6pt}
{\small Code and scripts: \url{https://github.com/Katherine-Deborah/AIVC-DA-perturbation}}

\end{document}
